\section{Introdução}

O câncer do colo do útero permanece entre as principais causas evitáveis de
morbidade e mortalidade por câncer em mulheres. Estimativas globais para 2020
indicam aproximadamente 604 mil novos casos e 342 mil mortes, com maior carga
em regiões que enfrentam restrições de acesso à vacinação, ao rastreamento e ao
tratamento oportuno \cite{singh2023global}. A citologia cervical continua sendo
um componente relevante do rastreamento, mas a leitura manual de lâminas é
trabalhosa, depende de especialistas e envolve alterações morfológicas que
podem ser sutis, especialmente nas categorias limítrofes.

Redes neurais profundas têm apresentado resultados promissores na análise de
imagens citológicas \cite{jiang2023systematic,xu2023cervical}. Entretanto,
revisões recentes mostram que a comparação entre métodos é dificultada por
diferenças de base, preparação da lâmina, tarefa, granularidade dos rótulos e
protocolo de validação \cite{fang2024systematic}. Um risco central ocorre
quando células recortadas da mesma imagem aparecem em treino e teste: elas
compartilham coloração, aquisição, fundo e artefatos, podendo inflar a
estimativa de generalização quando o particionamento não controla essa
dependência.

Além da discriminação, um sistema de apoio à triagem precisa tornar explícita a
relação entre falsos negativos e falsos positivos. Um limiar fixo de 0,5 pode
não ser adequado quando omitir uma célula anormal é mais grave do que
encaminhar uma célula normal para revisão. Probabilidades também exigem
calibração \cite{guo2017calibration}, e diretrizes como TRIPOD+AI reforçam a
necessidade de relatar partições, limiares, incerteza e uso pretendido
\cite{collins2024tripod}.

Este artigo estuda a classificação binária \emph{normal versus anormal} de
células da coleção CRIC Cervix, obtidas por citologia convencional. A tarefa é
formulada como triagem: ASC-US, LSIL, ASC-H, HSIL e SCC compõem a classe
anormal, mas os subtipos são preservados para análise posterior. Usa-se
ConvNeXt-Tiny pré-treinada \cite{liu2022convnext} para avaliar até que ponto um
classificador pode priorizar células suspeitas em imagens convencionais.

As contribuições são avaliar a ConvNeXt-Tiny na CRIC Cervix com separação por
imagem-fonte, analisar limiares operacionais de triagem, comparar
exploratoriamente com ResNet50 na mesma partição retida e interpretar os erros
pelas categorias Bethesda.
